\documentclass{article}
\usepackage[utf8]{inputenc}
\usepackage[T1]{fontenc}
\usepackage{geometry}
\geometry{
    a4paper,
    left=2.5cm,
    right=2.5cm,
    top=2.5cm,
    bottom=2.5cm
}
\usepackage{graphicx}
\usepackage{float}
\usepackage{hyperref}
\usepackage[english]{babel}
\usepackage{amsmath}
\usepackage{amssymb}
\usepackage{array}
\usepackage{booktabs}
\usepackage{caption}
\usepackage{subcaption}
\usepackage{enumitem}
\usepackage{fancyhdr}
\usepackage{titlesec}
\usepackage{url}
\usepackage{textcomp}
\usepackage{xcolor}
\usepackage{listings}
\usepackage{multirow}
\usepackage{longtable}
\usepackage{tabularx}
\usepackage{times}

% Set up page style
\pagestyle{fancy}
\fancyhf{}
\rfoot{\thepage}
\renewcommand{\headrulewidth}{0pt}

% Title formatting
\titleformat{\chapter}[display]{\normalfont\bfseries}{\chaptertitle\thechapter}{0pt}{\Huge}
\titlespacing*{\chapter}{0pt}{50pt}{40pt}

% Document info
\title{Template Document}
\author{Academic Support Agent}
\date{\today}

\begin{document}
\maketitle

\frame{\titlepage}
\section{Nombre de la secci n}
\begin{frame}
\frametitle{Nombre de la diapositiva}
Contenido de la diapositiva. Puede tener vi etas sin numerar
\begin{itemize}
    \item \emph{Se puede hacer un llamado de atenci n},
    \item O normalito,
     \begin{itemize}
         \item Se pueden colocar vi etas de segundo y n-nivel
     \end{itemize}
    \item Programaci n dirigida por eventos,
    \begin{enumerate}
        \item Numeradas tambi n
    \end{enumerate}
\end{itemize}
\begin{enumerate}
    \item O al rev s
    \begin{itemize}
        \item La numeraci n es en el primer nivel
    \end{itemize}
\end{enumerate}
\end{frame}

\section{Otra secci n - Esto aparece en el gestor de contenidos}
\begin{frame}
\frametitle{El nombre de la diapositiva}
Se permiten bloques de llamado de atenci n
\begin{alertblock}{No todo es blanco o negro}
A veces se puede llamar la atenci n.
\end{alertblock}
\begin{columns}
\column{0.5\textwidth}
Y qu  tal si se requiere texto en columnas?
\begin{enumerate}
    \item Desarrollo m s r pido,
\end{enumerate}
    
\column{0.5\textwidth}
Segunda columna\\
Comenzamos con las im genes:
\begin{figure}
    \centering
    \includegraphics[scale=0.2]{Logo_UniMinuto.png}
    \caption{El nombre de la figura}
    \label{fig:paraHacerLlamadosEnElCuerpo}
\end{figure}
\end{columns}
\end{frame}

\begin{frame}{Dentro de la misma secci n, otro t tulo}
Si se requieren v nculos:
\href{https://bit.ly/31Xqmaf}{Lenguajes}
\end{frame}

\subsection{Una subsecci n}
\begin{frame}{Siempre se nomenclan las diapositivas}
\href{https://www.youtube.com/watch?v=-SADbPS8UgA}{Arquitectura von Neumann}
\end{frame}

\section{Bibliograf a}
\begin{frame}{Tambi n hay espacio para las referencias}

\begin{thebibliography}{9}
\bibitem{Fernandez-Pena2011} 
Fern ndez-Pe a, Juan M. 
\textit{\href{https://www.uv.mx/personal/jfernandez/files/2010/07/8_Calidad.pdf}{Modelos de calidad de software}}. 2011

\bibitem{Pressman2010} 
Pressman, Roger S. 
\textit{\href{http://cotana.informatica.edu.bo/downloads/ld-Ingenieria.de.software.enfoque.practico.7ed.Pressman.PDF}{Ingenier a de Software: Un enfoque pr ctico 7a ed}}. Mc Graw-Hill, 2010

\bibitem{Callejas2017} 
Callejas-Cuervo, Mauro, Alarc n-Aldana, Andrea C. y  lvarez-Carre o, Ana M.
\textit{\href{hhttp://www.scielo.org.co/scielo.php?script=sci_abstract&pid=S1900-38032017000100236&lng=en&nrm=iso&tlng=es}{Modelos de calidad del software, un estado del arte}}. 2017

\bibitem{Fuertes2008}
Fuertes-Castro, Jos  L.
\textit{\href{http://www.iimv.org/iimv-wp-1-0/resources/uploads/2015/01/JoseLuisFuertesCalidaddelSoftware.pdf}{Calidad de Software}}. 2008
\end{thebibliography}
\end{frame}
\begin{frame}
\frametitle{Gracias}
\begin{figure}[h]
%\caption{Gracias}
\centering
\includegraphics[scale=0.6]{agradecimiento.jpg}
\end{figure}
\end{frame}
\end{document}